\subsection{Divisores relativos}
\label{subsection:IV.21.15}
\label{IV.21.15}

\begin{env}[21.15.1]
\label{IV.21.15.1}
Sean $S$ un preesquema y $f:X\to S$ un morfismo \emph{plano y localmente de
presentación finita}. Se ha definido en (20.6.1) el haz de anillos
$\mathscr M_{X/S}$ de los \emph{gérmenes de funciones meromorfas sobre $X$ relativas a
$S$}, subhaz de $\mathscr M_X$; es claro que la inyección canónica
$\mathcal O_X\to\mathscr M_X$ (20.1.4.1) aplica $\mathcal O_X$ en un
subhaz de $\mathscr M_{X/S}$, con el cual se lo identifica. Sea
$\mathscr M_{X/S}^*$ el haz (en grupos multiplicativos) de los gérmenes de secciones
invertibles de $\mathscr M_{X/S}$, que es, por tanto, un subhaz de $\mathscr M_X^*$
y contiene a $\mathcal O_X^*$ como subhaz.
\end{env}

\begin{definition}[21.15.2]
\label{IV.21.15.2}
Se llama \emph{haz de los divisores sobre $X$ relativos a $S$}, o
\emph{haz de los divisores sobre $X$ transversales a $f$}, y se denota por
$\mathscr{D}iv_{X/S}$ el haz cociente (de grupos conmutativos)
$\mathscr M_{X/S}^*/\mathcal O_X^*$; las secciones de este haz sobre $X$ se
denominan \emph{divisores sobre $X$ relativos a $S$}, o \emph{divisores sobre $X$
transversales a $f$}; forman un grupo conmutativo denotado por $\operatorname{Div}(X/S)$.
\end{definition}

Es claro que $\mathscr{D}iv_{X/S}$ se identifica canónicamente con un subhaz de
$\mathscr{D}iv_X$, y por tanto $\operatorname{Div}(X/S)$ con un subgrupo de
$\operatorname{Div}(X)$, que se denota también \emph{aditivamente}. Para
\oldpage[IV-4]{330}
todo abierto $U$ de $X$, se tiene
$\mathscr M_{X/S}^*|U=\mathscr M_{U/S}^*$, y por tanto
$\mathscr{D}iv_{X/S}|U=\mathscr{D}iv_{U/S}$, de modo que el haz
$\mathscr{D}iv_{X/S}$ es igual al prehaz
$U\mapsto\operatorname{Div}(U/S)$.

Puesto que $\mathscr M_{X/S}^*$ es un subhaz de $\mathscr M_X^*$, las
definiciones, notaciones y fórmulas relativas a los divisores de las secciones de
$\mathscr M_X^*$ sobre $X$ (21.1.3) se aplican sin cambio a las secciones
de $\mathscr M_{X/S}^*$ sobre $X$.

\begin{env}[21.15.3]
\label{IV.21.15.3}
La estructura de haz de grupos ordenados sobre $\mathscr{D}iv_X$ (21.1.6) induce
sobre el subhaz $\mathscr{D}iv_{X/S}$ una estructura de haz de grupos
ordenados, para la cual el haz de monoides de los gérmenes de secciones positivas es
$\mathscr{D}iv_X^+\cap\mathscr{D}iv_{X/S}$, que se denota por
$\mathscr{D}iv_{X/S}^+$. Se tiene
$\Gamma(X,\mathscr{D}iv_{X/S}^+)=\operatorname{Div}^+(X)\cap
\operatorname{Div}(X/S)$; se denota este submonoide de $\operatorname{Div}(X/S)$ por
$\operatorname{Div}^+(X/S)$, y está formado por los elementos $\geq0$ para una estructura
de grupo ordenado sobre $\operatorname{Div}(X/S)$. Resulta de (21.1.5.1) que
$\mathscr{D}iv_{X/S}^+$ es la imagen en $\mathscr{D}iv_{X/S}$ del subhaz de
monoides
\[
\label{IV.21.15.3.1}
  \mathcal O_{X/S}^{\mathrm{reg}}=\mathcal O_X\cap\mathscr M_{X/S}^*.
  \tag{21.15.3.1}
\]
\end{env}

Para todo abierto $U$ de $X$, $\Gamma(U,\mathcal O_{X/S}^{\mathrm{reg}})$ es
el conjunto de las secciones $t$ de $\mathcal O_X$ sobre $U$ tales que $t$ sea
regular y que $1/t$ pertenezca a $\Gamma(U,\mathscr M_{X/S})$, lo que significa,
con las notaciones de (20.6.1), que $t\in\mathscr T_{X/S}(U)$, de modo que el haz
$\mathcal O_{X/S}^{\mathrm{reg}}$ no es sino el haz denotado por $\mathscr T_{X/S}$
en (20.6.1). Se puede escribir, por tanto,
\[
\label{IV.21.15.3.2}
  \mathscr{D}iv_{X/S}^+=\mathcal O_{X/S}^{\mathrm{reg}}/\mathcal O_X^*
  \tag{21.15.3.2}
\]
salvo isomorfismo canónico.

\begin{env}[21.15.3.3]
\label{IV.21.15.3.3}
Sea $D\in\operatorname{Div}^+(X/S)$ y consideremos el subpreesquema cerrado $Y(D)$
de $X$ definido por el Ideal $\mathscr J_X(D)$ de $\mathcal O_X$ (21.6.5); en virtud de
lo anterior, para todo $x\in Y(D)$, existe un entorno abierto $U$ de $x$ en $X$
y una sección $t\in\mathscr T_{X/S}(U)$ tal que
$\mathscr J_X(D)|U=\mathcal O_U\cdot t$; puesto que $x\in Y(D)$, la imagen $t_{f(x)}$ de
$t$ en $\Gamma(U\cap X_{f(x)},\mathcal O_{X_{f(x)}})$ pertenece al ideal maximal
de $\mathcal O_{X_{f(x)},x}$ y además, por definición, para todo $s\in S$, la imagen
$t_s$ de $t$ en $\Gamma(U\cap X_s,\mathcal O_{X_s})$ es una sección
\emph{regular}. Se deduce, pues, de (11.3.8) y (19.2.4) que la inmersión canónica
$Y(D)\to X$ es \emph{transversalmente regular con respecto a $S$ y de
codimensión $1$} en el punto $x$. Siendo inmediata la recíproca, se ve que se pueden
\emph{identificar canónicamente} los divisores positivos sobre $X$ relativos a $S$ con los
subpreesquemas cerrados $Y$ de $X$ tales que la inyección canónica $Y\to X$ sea una
inmersión transversalmente regular con respecto a $S$ y de codimensión $1$. Haremos
habitualmente esta identificación.
\end{env}

\begin{proposition}[21.15.4]
\label{IV.21.15.4}
Sean $D$ un divisor sobre $X$ y $\mathscr J_X(D)$ el Ideal fraccionario invertible
correspondiente (21.2.5). Para que $D\in\operatorname{Div}(X/S)$, es necesario y suficiente
que, para todo $x\in X$, se tenga
$(\mathscr J_X(D)\cap\mathscr M_{X/S}^*)_x\ne0$ (o, equivalentemente, que
se tenga $\mathscr J_X(D)\subset\mathscr M_{X/S}$ y
$\mathscr J_X(D)^{-1}\subset\mathscr M_{X/S}$).
\end{proposition}

\begin{proof}
En efecto, decir que $D\in\operatorname{Div}(X/S)$ significa que, para todo $x\in X$,
existe un entorno abierto $U$ de $x$ y una sección
$f\in\Gamma(U,\mathscr M_{X/S}^*)$ tales que $D|U=\operatorname{div}_U(f)$; como
$\mathscr J_X(D)|U$ es el Ideal fraccionario invertible $\mathcal O_Uf$, la
proposición resulta inmediatamente.
\end{proof}

\oldpage[IV-4]{331}
\begin{proposition}[21.15.5]
\label{IV.21.15.5}
Sean $D$ un divisor sobre $X$, $\mathcal O_X(D)$ el Ideal fraccionario invertible y
$s_D$ la sección meromorfa regular de $\mathcal O_X(D)$ definidos canónicamente por
$D$ (21.2.8 y 21.2.9). Para que $D\in\operatorname{Div}(X/S)$, es necesario y suficiente
que $s_D\in\Gamma(X,\mathscr M_{X/S}(\mathcal O_X(D)))$.
\end{proposition}

\begin{proof}
En efecto, si $U$ es un abierto de $X$ tal que $D|U=\operatorname{div}_U(f)$, donde
$f\in\Gamma(U,\mathscr M_X^*)$, decir que
$s_D\in\Gamma(U,\mathscr M_{X/S}(\mathcal O_X(D)))$ significa, en virtud de las
definiciones (20.6.2), que $f^{-1}\in\Gamma(U,\mathscr M_{X/S})$, de donde resulta la
proposición.
\end{proof}

La interpretación de los divisores sobre $X$ mediante las clases
$\operatorname{cl}(\mathscr L,s)$ (21.2.11) permite, por tanto, interpretar los elementos de
$\operatorname{Div}(X/S)$ como los pares (salvo isomorfismo) $(\mathscr L,s)$
tales que $\mathscr L$ sea un $\mathcal O_X$-módulo invertible y que $s$ sea una
sección meromorfa de $\mathscr L$ sobre $X$, \emph{regular con respecto a
$S$} (20.6.5, (iii)).

\begin{proposition}[21.15.6]
\label{IV.21.15.6}
Sea $D$ un divisor sobre $X$ relativo a $S$, y supongamos que, para todo $x\in X$ tal
que $\operatorname{prof}(\mathcal O_{X_{f(x)},x})=1$, se tenga $D_x\geq0$
(resp. $D_x=0$). Entonces se tiene $D\geq0$ (resp. $D=0$).
\end{proposition}

\begin{proof}
Como en (21.1.8), podemos limitarnos al caso en que $D=\operatorname{div}(\varphi)$, donde
$\varphi$ es una función meromorfa regular relativa a $S$, y todo se reduce a ver
que, si $D_x\geq0$ en todo punto $x\in X$ tal que
$\operatorname{prof}(\mathcal O_{X_{f(x)},x})=1$, $\varphi$ está definida en todo
$X$. Pero esta hipótesis significa que, si $T=X-\operatorname{dom}(\varphi)$, se tiene
$\operatorname{prof}(\mathcal O_{X_{f(x)},x})\geq2$ para todo $x\in T$, y para
concluir basta aplicar (20.6.6).
\end{proof}

\begin{env}[21.15.7]
\label{IV.21.15.7}
Sea $X'$ un segundo $S$-preesquema plano y localmente de presentación finita sobre $S$, y
sea $g:X'\to X$ un $S$-morfismo. Si el $S$-morfismo $g$ es \emph{plano}, se sabe
(21.4.5) que la imagen inversa por $g$ de todo divisor sobre $X$ está definida; si
además $D\in\operatorname{Div}(X/S)$, resulta de la definición (21.15.2) y de
(20.6.8) que se tiene entonces $g^*(D)\in\operatorname{Div}(X'/S)$.
\end{env}

\begin{env}[21.15.8]
\label{IV.21.15.8}
Sea $X'$ un segundo $S$-preesquema plano y localmente de presentación finita sobre $S$, y
sea $g:X'\to X$ un $S$-morfismo \emph{finito y plano}. Obsérvese que $g$ es entonces
necesariamente de presentación finita (1.4.3, 1.4.6 y 1.6.3), por lo que
$g_*(\mathcal O_{X'})$ es un $\mathcal O_X$-módulo plano y de presentación finita, y
por tanto \emph{localmente libre} (2.1.12), es decir, $g$ es un morfismo
localmente libre (18.2.7); para todo $s\in S$, el morfismo correspondiente
$g_s:X'_s\to X_s$ es entonces también finito y localmente libre. Se deduce de (21.5.2)
y (20.6.1) que, para todo abierto $U$ de $X$ y toda sección
$t'\in\Gamma(g^{-1}(U),\mathscr T_{X'/S})$, la norma $N_{X'/X}(t')$ pertenece a
$\Gamma(U,\mathscr T_{X/S})$; el razonamiento de (21.5.3) demuestra luego que, para todo
$\mathcal O_{X'}$-módulo invertible $\mathscr L'$ y toda sección meromorfa $u'$ de
$\mathscr L'$ sobre $X'$, \emph{regular con respecto a $S$}, la norma
$N_{X'/X}(u')$ es una sección meromorfa de
$\mathscr L=N_{X'/X}(\mathscr L')$ sobre $X$, regular con respecto a $S$.
La interpretación de los divisores relativos a $S$ dada en (21.15.5) y la definición de
la imagen directa de un divisor (21.5.5) demuestran entonces que \emph{para todo divisor}
$D'\in\operatorname{Div}(X'/S)$, se tiene $g_*(D')\in\operatorname{Div}(X/S)$.
\end{env}

\begin{env}[21.15.9]
\label{IV.21.15.9}
Consideremos, por último, un morfismo cualquiera $S'\to S$, y (bajo las hipótesis de
(21.15.1)) pongamos $X'=X\times_S S'$, que es plano y localmente de presentación finita
sobre $S'$; si $p:X'\to X$ es la proyección canónica, se ha visto (20.6.9) que se tiene un
homomorfismo canónico
$p^*(\mathscr M_{X/S})\to\mathscr M_{X'/S'}$, que transforma evidentemente toda
sección de $\mathscr M_{X/S}$ sobre un abierto $U$, \emph{regular con respecto
a $S$}, en una sección de $\mathscr M_{X'/S'}$ sobre $p^{-1}(U)$, regular
con respecto a $S'$\footnote{La transcripción francesa dice aquí «relativa a $S$»; tras el cambio de base, la sección pertenece a $\mathscr M_{X'/S'}$ y es regular con respecto a $S'$, como confirma el testigo inglés mantenido.} (20.6.5, (iii)); se concluye entonces de la
\oldpage[IV-4]{332}
definición (21.15.2) y de la exactitud por la derecha del funtor $p^*$ que
el homomorfismo anterior define, pasando a los cocientes, un homomorfismo canónico
\[
\label{IV.21.15.9.1}
  \operatorname{Div}(X/S)\longrightarrow\operatorname{Div}(X'/S')
  \tag{21.15.9.1}
\]
que transforma evidentemente los elementos de $\operatorname{Div}^+(X/S)$ en elementos de
$\operatorname{Div}^+(X'/S')$. Se pone también
$\operatorname{Div}(X'/S')=\operatorname{Div}_{X/S}(S')$ (resp.
$\operatorname{Div}^+(X'/S')=\operatorname{Div}_{X/S}^+(S')$), y se ve inmediatamente
que se han definido así dos funtores contravariantes
\[
  \operatorname{Div}_{X/S}:\mathbf{Sch}_{/S}\longrightarrow\mathbf{Ab},\qquad
  \operatorname{Div}_{X/S}^+:\mathbf{Sch}_{/S}\longrightarrow\mathbf{Ens}
\]
de la categoría de los $S$-preesquemas en la de los grupos conmutativos (resp. de los
conjuntos). Se verán más adelante (cap.~VI) casos importantes en los que el funtor
$\operatorname{Div}_{X/S}^+$ es \emph{representable} $(\mathbf 0_{\mathrm{III}},
8.1.8)$.

Para todo divisor $D\in\operatorname{Div}(X/S)$, la imagen de $D$ por
el homomorfismo (21.15.9.1) no es otra que la \emph{imagen inversa}
$p^*(D)$ (en el sentido de (21.4.2)): la existencia de esta imagen inversa y la afirmación
anterior son, en efecto, consecuencias inmediatas de (20.6.5, (iii)) y (20.6.9).

Los elementos de $\operatorname{Div}(X'/S')$ se llaman a menudo, por abuso de lenguaje,
«\emph{familias de divisores sobre $X$ relativos a $S$, parametrizadas por el
$S$-preesquema $S'$}»; esta terminología se utiliza sobre todo cuando se trata de
divisores \emph{positivos}.
\end{env}
